\section{Code data：The Stack 与 Stack v2}

The Stack 从 GitHub Archive 获取 repo names，git clone 大量 repositories，只保留 permissively licensed 代码，用 go-license-detector 检测 license，并用 MinHash/Jaccard 去重。Stack v2 扩展到 issues、comments、PRs、Software Heritage repos 和 docs crawls。

\begin{figure}[H]
\centering
\includegraphics[width=0.38\textwidth]{stackv2-pr1.png}
\hfill
\includegraphics[width=0.52\textwidth]{stackv2-pr2.png}
\caption{Stack v2 pull request 数据：把结构化 PR 对象线性化成 token sequence，并加入 diff 周边上下文。}
\figsource{CS336 Lecture 13 官方源引用图。}
\end{figure}

\begin{knowledgebox}{读图：为什么 PR metadata 有价值}
代码文件教模型语法和 API，PR/issue/comment 教模型软件工程过程：bug report、review、patch rationale、discussion。Agentic coding models 很可能受益于这种过程数据，但也要处理 PII、bot activity、malware 和 license。
\end{knowledgebox}

\begin{importantbox}{代码数据的特殊性}
代码不是普通文本：它有许可证、依赖、可执行语义、测试、提交历史和安全风险。高质量代码数据不只是 .py/.js 文件，还包括 README、issue、PR、review、diff context、CI log。面向 coding agents 的数据尤其需要过程信息。
\end{importantbox}

